\begin{resumo}[Abstract]
 \begin{otherlanguage*}{english}
This work presents SIGCE, a web system for managing internship hiring for students of the Faculty of Sciences and Technologies in Engineering (FCTE) at the University of Brasília. The process covers the internship request, advisor assignment, and the path of the internship agreement through review and the advisor's signature. Development was applied and organized in incremental deliveries. Requirements were elicited in a brainstorming session with students who had already completed internships and with the professor who runs the process. The system is a single application, with a Next.js frontend, a Spring Boot backend, and PostgreSQL, structured in layers according to Clean Architecture. Access uses the Brazilian federal single sign-on, Login Único. Advisor assignment reuses the radius rule already used by the internship commission's Telegram bot. The agreement uploaded by the student is checked by a generative model, which reports pending items and does not replace the professor's review. The student, the company supervisor, and the advisor each have a distinct signature step. Calling the Gov.br electronic-signature API still depends on Login Único being approved for production. The result is a platform that concentrates process tracking at FCTE and states both the limit of the automated review and the signature integration that remains pending.

   \vspace{\onelineskip}
 
   \noindent 
   \textbf{Key-words}: internship management. single sign-on. internship agreement. FCTE.
 \end{otherlanguage*}
\end{resumo}
